\section{Požiadavky na systém}

Na základe konzultácií s vedúcim projektu a analýzy potrieb cieľových používateľov bol zostavený zoznam funkčných a nefunkčných požiadaviek na systém. Požiadavky sú rozdelené do niekoľkých logických kategórií, ktoré odrážajú hlavné oblasti funkcionality platformy. Pri každej požiadavke je uvedený stav implementácie v rámci projektu.

\subsection{Funkčné požiadavky}

\subsubsection{Spracovanie šifrovacích dokumentov}

Jadrom systému je poloautomatické spracovanie historických šifrovacích dokumentov realizované prostredníctvom šesťkrokového pipeline rozhrania.

\begin{itemize}
    \item[\checkmark] Systém umožňuje nahratie dokumentu (obrázok, sken) s podporou drag\,\&\,drop a automaticky validuje formát dokumentu, pričom neplatné vstupy (obrázky bez textu) sú odmietnuté pomocou OCR kontroly.
    \item[\checkmark] Používateľ má možnosť upraviť parametre predspracovania obrazu (natočenie, jas, kontrast, automatická korekcia naklonenia, binarizácia, odstránenie šumu, korekcia tieňovania, odstránenie škvŕn) s okamžitým náhľadom výsledku.
    \item[\checkmark] Systém automaticky deteguje oblasti záujmu (sekcie, páry a elementy) v šifrovacích kľúčoch pomocou modelu YOLOv11 a zobrazí výsledky vo forme interaktívnych ohraničujúcich rámcov.
    \item[\checkmark] Používateľ má možnosť manuálne upraviť bounding boxy po detekcii, vytvárať nové, presúvať, meniť ich veľkosť, prejmenovávať triedy a mazať ich.
    \item[\checkmark] Systém aplikuje OCR (Tesseract) na detegované elementy a priraďuje im transkripciu textu.
    \item[\checkmark] Používateľ môže manuálne reklasifikovať elementy na \textit{plain text}, \textit{cipher text} alebo \textit{symbol} a opraviť OCR transkripciu.
    \item[\checkmark] Systém navrhuje priradenie symbolov z databázy ku každému \textit{cipher text} elementu (funkcia momentálne implementovaná ako mock; integrácia so skutočnou databázou symbolov je plánovaná).
    \item[\checkmark] Výsledky spracovania sú exportovateľné vo formáte JSON s hierarchickou štruktúrou (sekcie → páry → elementy).
    \item[$\circ$] Mapovanie symbolov do reálnej centralizovanej databázy a rekonštrukcia substitučných tabuliek sú predmetom ďalšieho vývoja.
    \item[$\circ$] Dešifrovanie textu s využitím rekonstruovaného kľúča je plánované ako budúca funkcionalita.
\end{itemize}

\subsubsection{Používateľský profil a autentifikácia}

\begin{itemize}
    \item[\checkmark] Systém umožňuje registráciu nového používateľa s overením e-mailovej adresy.
    \item[\checkmark] Používateľ sa môže registrovať a prihlasovať pomocou Google účtu (OAuth 2.0).
    \item[\checkmark] Prihlásenie a odhlásenie sú plne implementované; autentifikačný stav je zachovaný medzi reláciami pomocou JWT tokenu uloženého v \texttt{localStorage}.
    \item[\checkmark] Používateľ môže upraviť svoje používateľské meno a heslo v sekcii profilu.
    \item[\checkmark] Systém zaznamenáva aktivitu používateľa (nahranie, úprava, zdieľanie, mazanie dokumentov, prihlásenie, odhlásenie) do denníka aktivít.
\end{itemize}

\subsubsection{Správa dokumentov}

\begin{itemize}
    \item[\checkmark] Systém zobrazuje všetky dokumenty prihláseného používateľa v prehľadnom mriežkovom rozložení.
    \item[\checkmark] Používateľ môže dokumenty vyhľadávať podľa názvu.
    \item[\checkmark] Dokumenty je možné otvoriť priamo v pipeline rozhraní pre ďalšie spracovanie.
    \item[\checkmark] Systém podporuje prepínanie dokumentov medzi režimami \texttt{private}, \texttt{public} a \texttt{shared}.
    \item[\checkmark] Dokumenty je možné zdieľať s konkrétnymi používateľmi; zdieľané dokumenty sú prístupné v sekcii \textit{Shared with me}.
    \item[\checkmark] Verejné dokumenty sú dostupné všetkým prihláseným používateľom v sekcii \textit{Public documents}.
    \item[$\circ$] Filtrovanie dokumentov podľa typu, jazyka, krajiny pôvodu alebo dátumu je plánované ako budúca funkcionalita.
    \item[$\circ$] Nastavenie rolí prístupu (čítanie / zápis) pre zdieľanie s inými používateľmi je predmetom ďalšieho vývoja.
\end{itemize}

\subsubsection{Administrácia}

\begin{itemize}
    \item[\checkmark] Administrátor má prístup k prehľadu všetkých používateľov a môže spravovať ich účty.
    \item[\checkmark] Administrátor má prístup k prehľadu všetkých dokumentov v systéme.
    \item[$\circ$] Správa symbolovej databázy (normalizácia, zlúčenie duplicít, odstránenie chýb) je plánovaná ako budúca funkcionalita po zavedení reálnej databázy symbolov.
\end{itemize}

\subsubsection{Gamifikácia}

Pôvodne plánovaný systém bodov, kreditov a rebríčka používateľov bol v priebehu projektu vyhodnotený ako mimo rozsahu aktuálnej implementácie a jeho realizácia bola odložená na neskorší vývoj platformy.

\subsection{Nefunkčné požiadavky}

\begin{itemize}
    \item[\checkmark] Systém je dostupný cez webové rozhranie na doméne \texttt{xmolnarf1.com}.
    \item[\checkmark] Bezpečnosť je zabezpečená bcrypt hashovaním hesiel, JWT autentifikáciou a kontrolou prístupov k dokumentom.
    \item[\checkmark] Komunikácia so serverom prebieha šifrovane cez HTTPS (SSL certifikát spravovaný Cloudflare).
    \item[\checkmark] Používateľské rozhranie je intuitívne a responzívne vrátane podpory pre mobilné zariadenia.
    \item[\checkmark] Architektúra systému je modulárna – AI modul, backend a frontend sú nasadené ako nezávislé Docker kontajnery, čo umožňuje jednoduché rozšírenie o nové AI moduly.
\end{itemize}
